\documentclass[12pt,a4paper]{article}

% ── Packages ──────────────────────────────────────────────────────────────────
\usepackage[margin=2.5cm, headheight=14pt]{geometry}
\usepackage{amsmath, amssymb, amsthm}
\usepackage{mathtools}
\usepackage{booktabs}
\usepackage{array}
\usepackage{longtable}
\usepackage{multirow}
\usepackage{graphicx}
\usepackage{xcolor}
\usepackage{hyperref}
\usepackage{enumitem}
\usepackage{fancyhdr}
\usepackage{titlesec}
\usepackage{parskip}
\usepackage{float}
\usepackage{caption}
\usepackage{subcaption}
\usepackage{listings}
\usepackage{tikz}
\usepackage{pgfplots}
\pgfplotsset{compat=1.18}

% ── Formatting ────────────────────────────────────────────────────────────────
\hypersetup{
  colorlinks=true,
  linkcolor=blue!70!black,
  citecolor=blue!70!black,
  urlcolor=blue!70!black,
  pdftitle={Bank Balance Sheet Optimization},
  pdfauthor={Apex Global Bank — Quant Research}
}

\pagestyle{fancy}
\fancyhf{}
\lhead{\small Apex Global Bank — Proprietary \& Confidential}
\rhead{\small Balance Sheet Optimization}
\cfoot{\small \thepage}
\renewcommand{\headrulewidth}{0.4pt}

\titleformat{\section}{\large\bfseries}{}{0em}{\thesection\quad}
\titleformat{\subsection}{\normalsize\bfseries}{}{0em}{\thesubsection\quad}

% ── Custom commands ────────────────────────────────────────────────────────────
\newcommand{\basis}{\text{bp}}
\newcommand{\E}{\mathbb{E}}
\newcommand{\R}{\mathbb{R}}
\DeclareMathOperator*{\argmax}{arg\,max}
\DeclareMathOperator*{\argmin}{arg\,min}

\newtheorem{definition}{Definition}
\newtheorem{proposition}{Proposition}
\newtheorem{remark}{Remark}

% ── Document ──────────────────────────────────────────────────────────────────
\begin{document}

\begin{titlepage}
  \centering
  \vspace*{2cm}
  {\Large\bfseries Apex Global Bank\\[0.4em]}
  {\large\bfseries Quantitative Research \& Treasury\par}
  \vspace{1.5cm}
  {\LARGE\bfseries Bank Balance Sheet Optimization\\[0.5em]}
  {\large Theory, Mathematics, Regulatory Framework,\\
   Technology, and Revenue Generation\par}
  \vspace{1cm}
  \rule{\textwidth}{0.4pt}
  \vspace{0.5cm}
  {\large Version 1.0 \quad|\quad April 2026\par}
  \vspace{0.3cm}
  {\normalsize Prepared by: Dr.\ Yuki Tanaka (Head of Quant Research)\\
   Reviewed by: Dr.\ Priya Nair (CRO) \& CFO Office\par}
  \vspace{0.5cm}
  \rule{\textwidth}{0.4pt}
  \vspace{1cm}
  \begin{minipage}{0.85\textwidth}
    \small\itshape
    This document is proprietary to Apex Global Bank.  It describes
    the quantitative framework, regulatory constraints, computational
    infrastructure, and revenue-generation mechanisms underlying the
    bank's balance sheet optimization program.  Sections 3--8 contain
    material non-public information.  Distribution restricted to
    authorized personnel.
  \end{minipage}
  \vfill
\end{titlepage}

\tableofcontents
\newpage

% ══════════════════════════════════════════════════════════════════════════════
\section{Introduction and Strategic Context}
% ══════════════════════════════════════════════════════════════════════════════

\subsection{What Is Balance Sheet Optimization?}

A bank's balance sheet is fundamentally a portfolio of financial claims.  On
the asset side: loans, securities, derivatives, and cash.  On the liability
side: deposits, wholesale funding, issued debt, and equity.  \emph{Balance
sheet optimization} is the discipline of constructing and managing this
portfolio to maximize risk-adjusted return subject to regulatory constraints,
liquidity requirements, and the bank's own risk appetite.

Unlike trading book optimization, which operates on a horizon of hours to
days, balance sheet optimization operates on horizons of months to years.
It is a strategic planning function that touches every product line:
commercial lending, retail deposits, fixed-income securities portfolios,
derivatives books, repo/securities finance, and off-balance-sheet vehicles.

At the highest level, the problem is:

\begin{equation}
  \max_{\mathbf{a},\,\mathbf{l}} \quad \text{RORWA}(\mathbf{a}, \mathbf{l})
  \quad \text{subject to} \quad
  \mathbf{g}(\mathbf{a}, \mathbf{l}) \leq \mathbf{0}
\end{equation}

where $\mathbf{a}$ is the asset allocation vector, $\mathbf{l}$ is the
liability structure vector, RORWA is the return on risk-weighted assets, and
$\mathbf{g}$ encodes the regulatory and risk constraints (capital ratios,
liquidity ratios, concentration limits, stress test floors).

In practice this is far more complex than a single optimization: the
constraints are multi-dimensional, some are non-linear, many are path-dependent
over time, and the objective itself is a function of models that embed their
own uncertainty.

\subsection{Why It Matters: Revenue at Stake}

At a JPMorgan-scale institution (\$3.9 trillion balance sheet), a 5 basis
point improvement in return on assets (ROA) is worth approximately
\$2 billion per year.  The levers available through balance sheet optimization
include:

\begin{itemize}[noitemsep]
  \item \textbf{Capital efficiency}: replacing high-RWA assets with
        equivalent-yield lower-RWA assets releases CET1 that can be redeployed
        or returned to shareholders via buybacks.
  \item \textbf{Funding cost reduction}: optimizing the liability mix
        (longer-term funding when the curve is inverted; shorter-term when
        normally sloped) reduces the weighted average cost of funds.
  \item \textbf{NII maximization}: positioning the balance sheet's
        asset-liability duration gap to benefit from expected rate moves.
  \item \textbf{Regulatory arbitrage} (within the rules): choosing product
        structures that achieve the same economic exposure with lower capital
        consumption (e.g., guarantees vs.\ direct loans, cleared vs.\
        bilateral derivatives).
  \item \textbf{Liquidity buffer efficiency}: minimizing the drag from
        mandatory HQLA holdings while maintaining LCR/NSFR compliance.
\end{itemize}

\subsection{The Regulatory Stack}

Balance sheet optimization operates inside a dense regulatory framework.  The
binding constraints are:

\begin{table}[H]
  \centering
  \caption{Key Regulatory Constraints on the Bank Balance Sheet}
  \begin{tabular}{llll}
    \toprule
    Regulation & Metric & Minimum & Apex Current \\
    \midrule
    Basel III / 12 CFR 3 & CET1 ratio & 4.5\% (+ CCB 2.5\%) & 13.0\% \\
    Basel III / 12 CFR 3 & Tier 1 ratio & 6.0\% (+ CCB) & 14.5\% \\
    Basel III / 12 CFR 3 & Total Capital ratio & 8.0\% (+ CCB) & 16.2\% \\
    Basel III / 12 CFR 3 & Leverage ratio & 3.0\% (SLR) & 7.8\% \\
    Basel III / 12 CFR 249 & LCR & 100\% & 118\% \\
    Basel III / 12 CFR 249 & NSFR & 100\% & 112\% \\
    Dodd-Frank / 12 CFR 252 & SCB / DFAST CET1 floor & 4.5\% + SCB (regulatory) & 8.7\% trough \\
    Basel IV / FRTB & FRTB-SA floor & Binding from Jan 2026 (US) & --- \\
    Basel III / 12 CFR 3 & G-SIB surcharge (CET1) & +1.0\%--3.5\% additive & +2.5\% (Bucket 3) \\
    Basel III / 12 CFR 3 & CCyB (countercyclical) & 0\%--2.5\% (jurisdiction) & 0\% (US current) \\
    Concentration rules & Single-name limit & 25\% of CET1 & Monitored \\
    \bottomrule
  \end{tabular}
\end{table}

Each of these constraints reduces the feasible set and potentially creates
binding limits that shape the optimal balance sheet structure.

% ══════════════════════════════════════════════════════════════════════════════
\section{The Optimization Problem: Mathematical Formulation}
% ══════════════════════════════════════════════════════════════════════════════

\subsection{State Variables and Decision Variables}

Let $t = 0, 1, \ldots, T$ index quarterly periods.  The state of the balance
sheet at time $t$ is described by:

\begin{itemize}[noitemsep]
  \item $A_k(t)$: balance in asset class $k$ (loans, AFS securities, HTM
        securities, derivatives, cash, other), for $k = 1,\ldots,K_A$
  \item $L_j(t)$: balance in liability class $j$ (demand deposits, time
        deposits, FHLB advances, senior unsecured debt, repo, equity),
        for $j = 1,\ldots,K_L$
  \item $r_k^A(t)$: yield on asset class $k$ at time $t$
  \item $r_j^L(t)$: cost of liability class $j$ at time $t$
  \item $\rho_k(t)$: risk weight for asset class $k$ (SA or IMA)
  \item $q_t$: macroeconomic state vector (GDP growth, unemployment, rates)
\end{itemize}

The decision variables are the \emph{incremental} changes:

\begin{equation}
  \Delta A_k(t), \quad \Delta L_j(t), \quad \forall\, k, j, t
\end{equation}

subject to the balance sheet identity:

\begin{equation}
  \sum_k A_k(t) = \sum_j L_j(t), \qquad \forall\, t
  \label{eq:bsi}
\end{equation}

\subsection{Objective Function}

The primary objective is to maximize the net present value of risk-adjusted
returns over the planning horizon $[0, T]$:

\begin{equation}
  \max_{\{\Delta A_k(t),\, \Delta L_j(t)\}}
  \sum_{t=0}^{T} \delta^t \left[
    \text{NII}(t) + \text{Fee}(t) - \text{Credit\_Loss}(t)
    - \text{OpRisk}(t) - \text{FTP\_Cost}(t)
  \right]
\end{equation}

where $\delta = 1/(1+h)$ is the discount factor (hurdle rate $h$ is the
bank's cost of equity, typically CAPM-derived: $h = r_f + \beta (r_m - r_f)$).

Equivalently, using the RORWA metric preferred by bank management:

\begin{equation}
  \text{RORWA}(t) = \frac{\text{Net\_Income\_after\_tax}(t)}{\text{RWA}(t)}
\end{equation}

\subsection{Net Interest Income}

\begin{equation}
  \text{NII}(t) = \sum_k r_k^A(t) \cdot A_k(t) - \sum_j r_j^L(t) \cdot L_j(t)
\end{equation}

The net interest margin (NIM) is:

\begin{equation}
  \text{NIM}(t) = \frac{\text{NII}(t)}{\text{Average\_Earning\_Assets}(t)}
\end{equation}

\subsection{Capital Constraints}

The CET1 constraint:

\begin{equation}
  \frac{\text{CET1}(t)}{\text{RWA}(t)} \geq \kappa_{\text{CET1}}, \qquad \forall\, t
  \label{eq:cet1}
\end{equation}

where $\kappa_{\text{CET1}} = 7\%$ (4.5\% minimum + 2.5\% CCB) in normal conditions
and $\kappa_{\text{CET1}}^{\text{stress}} = 6.5\%$ under severely adverse DFAST.

Risk-weighted assets:

\begin{equation}
  \text{RWA}(t) = \sum_k \rho_k(t) \cdot A_k(t)
  \label{eq:rwa}
\end{equation}

The leverage ratio constraint:

\begin{equation}
  \frac{\text{Tier1}(t)}{\text{Total\_Leverage\_Exposure}(t)} \geq \kappa_{\text{SLR}}
  \label{eq:slr}
\end{equation}

where Total Leverage Exposure = total assets + off-balance-sheet items + derivative
add-ons.  The SLR (Supplementary Leverage Ratio) is often the binding constraint
for large banks with low-RWA assets (Treasuries, Agency MBS).

\subsection{Liquidity Constraints}

\textbf{LCR (Liquidity Coverage Ratio):}

\begin{equation}
  \text{LCR}(t) = \frac{\text{HQLA}(t)}{\text{Net\_Cash\_Outflows}_{30d}(t)} \geq 1.0
  \label{eq:lcr}
\end{equation}

where HQLA (High-Quality Liquid Assets) = Level 1 assets (cash, central bank
reserves, Treasuries) at 100\% + Level 2A assets (Agency debt, covered bonds)
at 85\% + Level 2B assets (non-financial corporate bonds, equities) at 50\%.

Net cash outflows are computed from stressed runoff rates applied to each
liability category:

\begin{equation}
  \text{Net\_Cash\_Outflows}_{30d}(t) = \sum_j \alpha_j \cdot L_j(t) - \text{Cash\_Inflows}(t)
\end{equation}

where $\alpha_j$ is the 30-day stressed runoff rate for liability class $j$
(e.g., $\alpha_{\text{retail deposits}} = 3\%$, $\alpha_{\text{wholesale unsecured}} = 25\%$,
$\alpha_{\text{secured funding}} = 0\%\text{ to } 25\%$ depending on collateral quality).

\textbf{NSFR (Net Stable Funding Ratio):}

\begin{equation}
  \text{NSFR}(t) = \frac{\text{ASF}(t)}{\text{RSF}(t)} \geq 1.0
  \label{eq:nsfr}
\end{equation}

Available Stable Funding (ASF):
\begin{equation}
  \text{ASF}(t) = \sum_j \beta_j^{\text{ASF}} \cdot L_j(t)
\end{equation}

Required Stable Funding (RSF):
\begin{equation}
  \text{RSF}(t) = \sum_k \beta_k^{\text{RSF}} \cdot A_k(t)
\end{equation}

where $\beta^{\text{ASF}}$ and $\beta^{\text{RSF}}$ are Basel-prescribed factors.
High-RSF assets (e.g., loans > 1 year, RSF = 65--85\%) paired with low-ASF
short-term funding creates NSFR pressure.

\subsection{DFAST / Stress Capital Buffer Constraint}

The stress test constraint ensures the CET1 ratio survives the 9-quarter
projection even under the severely adverse scenario:

\begin{equation}
  \min_{t \in [1,9]} \text{CET1\_ratio}^{\text{sev\_adv}}(t) \geq \kappa_{\text{floor}}
  \label{eq:dfast}
\end{equation}

where $\kappa_{\text{floor}} = 6.5\%$ (internal risk appetite) or 4.5\% (regulatory minimum).

\begin{remark}[Stress Capital Buffer]
Since 2020 (12 CFR 252.70), the Federal Reserve replaced the old DFAST-based
capital action restrictions with the \emph{Stress Capital Buffer} (SCB).  The
SCB equals the maximum decline in the CET1 ratio during the severely adverse
scenario (subject to a 2.5\% floor), and is added directly to the firm's
minimum CET1 requirement.  A bank that projects a 3.5\% CET1 drawdown in
the severe scenario faces a minimum CET1 requirement of $4.5\% + 3.5\% = 8.0\%$
(before the G-SIB surcharge and CCyB).  Balance sheet optimization must
therefore target a CET1 level sufficient to pass the SCB test, not merely a
static 7\% floor.
\end{remark}

This constraint is forward-looking and non-linear: it couples the current balance
sheet composition to the modelled stressed loss path over 9 quarters.

\subsection{Concentration Constraints}

Single-name concentration:
\begin{equation}
  A_{\text{counterparty-}i}(t) \leq 0.25 \cdot \text{CET1}(t), \qquad \forall\, i
  \label{eq:conc}
\end{equation}

Sector concentration:
\begin{equation}
  \sum_{i \in \text{sector-}s} A_i(t) \leq 0.25 \cdot \text{Total\_Assets}(t), \qquad \forall\, s
\end{equation}

\subsection{The Full Constrained Problem}

Combining constraints \eqref{eq:bsi}--\eqref{eq:conc}:

\begin{equation}
  \begin{aligned}
    \max_{\Delta A, \Delta L} \quad &
      \sum_{t=0}^{T} \delta^t \left[ \text{NII}(t) - \text{CreditLoss}(t)
      - \text{OpRisk}(t) \right] \\
    \text{s.t.} \quad
    & \eqref{eq:bsi}: \text{Balance sheet identity} \\
    & \eqref{eq:cet1}: \text{CET1 ratio} \geq 7\% \\
    & \eqref{eq:slr}: \text{SLR} \geq 3\% \\
    & \eqref{eq:lcr}: \text{LCR} \geq 100\% \\
    & \eqref{eq:nsfr}: \text{NSFR} \geq 100\% \\
    & \eqref{eq:dfast}: \text{DFAST CET1 trough} \geq 6.5\% \\
    & \eqref{eq:conc}: \text{Concentration limits} \\
    & A_k(t) \geq 0, \quad L_j(t) \geq 0 \quad \forall\, k, j, t
  \end{aligned}
  \label{eq:fullopt}
\end{equation}

This is a multi-period nonlinear program (NLP) with integer components
(lending decisions), stochastic elements (credit losses depend on macro
realization), and regulatory feedback (RWA changes with asset mix).  In
practice it is decomposed into subproblems: capital optimization, liquidity
optimization, NII/duration management, and pricing.

% ══════════════════════════════════════════════════════════════════════════════
\section{Capital Optimization}
% ══════════════════════════════════════════════════════════════════════════════

\subsection{RWA Minimization and the Capital Efficiency Frontier}

The \emph{capital efficiency frontier} maps every unit of expected return
against the required capital.  The Sharpe-style metric is:

\begin{equation}
  \text{RAROC}(k) = \frac{r_k^A - r_k^{\text{FTP}} - \text{EL}_k}
                         {\text{EC}_k}
\end{equation}

where:
\begin{itemize}[noitemsep]
  \item $r_k^A$: gross yield on asset class $k$
  \item $r_k^{\text{FTP}}$: matched-maturity FTP charge (Section~6)
  \item $\text{EL}_k = \text{PD}_k \cdot \text{LGD}_k$: expected loss rate
  \item $\text{EC}_k$: economic capital allocation per unit (see below)
\end{itemize}

A business line is value-creating if $\text{RAROC}(k) > h$ (hurdle rate).  The
optimization re-allocates the balance sheet toward high-RAROC activities.

\subsection{Economic Capital Allocation}

Regulatory capital (RWA-based) and economic capital (model-based) can
diverge.  Economic capital is typically computed using a portfolio credit VaR
at a 99.9\% confidence level (consistent with BBB-rated survival probability):

\begin{equation}
  \text{EC}_{\text{portfolio}} = \text{VaR}_{99.9\%}(\text{CreditLoss}_{1yr}) - \text{EL}
\end{equation}

Per-obligor capital is allocated using the marginal contribution:

\begin{equation}
  \text{EC}_i = \frac{\partial \text{EC}_{\text{portfolio}}}{\partial w_i}
              = \text{Corr}(L_i, L_{\text{portfolio}}) \cdot \sigma_{L_i} \cdot z_{99.9\%}
\end{equation}

where $L_i$ is the loss random variable for obligor $i$.  This ``marginal
EC'' framework means that adding a low-correlation obligor consumes less
capital than an equally risky but highly correlated one.

\subsection{SA vs.\ IMA Capital Arbitrage}

Under Basel III, banks with IMA approval compute market risk capital using
internal VaR models (lower capital for diversified portfolios).  Under FRTB,
all banks report SA capital as a floor.

The capital ratio under SA vs.\ IMA:

\begin{equation}
  \text{RWA}^{\text{SA}}(k) = \rho_k^{\text{SA}} \cdot \text{EAD}_k
\end{equation}

\begin{equation}
  \text{RWA}^{\text{IMA}}(k) = 12.5 \cdot \left[ \max\left(\text{VaR}_t, \bar{k} \cdot \overline{\text{VaR}}_{60d}\right)
                                         + \max\left(\text{SVaR}_t, \bar{k} \cdot \overline{\text{SVaR}}_{60d}\right) \right]
\end{equation}

For a well-diversified trading book, $\text{RWA}^{\text{IMA}} < \text{RWA}^{\text{SA}}$ by
30--60\%, representing a significant capital saving that justifies the cost of
IMA model development and validation.

\subsection{Capital Buffer Tiering}

Optimal capital management requires maintaining a buffer above the regulatory
minimum.  The buffer serves multiple purposes:

\begin{equation}
  \text{CET1}_{\text{target}} = \kappa_{\text{min}} + B_{\text{CCB}} + B_{\text{G-SIB}}
    + B_{\text{CCyB}} + B_{\text{SCB}} + B_{\text{management}}
\end{equation}

\begin{itemize}[noitemsep]
  \item $\kappa_{\text{min}} = 4.5\%$: Basel III minimum
  \item $B_{\text{CCB}} = 2.5\%$: Capital Conservation Buffer (restrictions on
        distributions trigger below this level)
  \item $B_{\text{G-SIB}} = 1.0\text{--}3.5\%$: G-SIB surcharge assessed annually
        (Apex: $+2.5\%$ at Bucket 3 based on systemic importance scores)
  \item $B_{\text{CCyB}} = 0\text{--}2.5\%$: Countercyclical Capital Buffer, set
        by each national regulator; currently 0\% in the US but active in
        several European and Asian jurisdictions; must be monitored for
        cross-border portfolios
  \item $B_{\text{SCB}}$: Stress Capital Buffer (max CET1 drawdown in DFAST
        severely adverse, floor 2.5\%); at Apex $\approx 3.5\%$ based on most
        recent stress test results
  \item $B_{\text{management}} = 1.0\text{--}1.5\%$: management buffer for
        strategic flexibility, acquisitions, unexpected losses
\end{itemize}

\begin{remark}[TLAC / MREL]
G-SIBs face a further constraint on liability structure: the \emph{Total
Loss-Absorbing Capacity} (TLAC) requirement (12 CFR 252.61--252.65 in the US;
MREL in the EU).  TLAC requires a minimum of $18\%$ of RWA (or $6.75\%$ of
leverage exposure) in loss-absorbing instruments (CET1, AT1, Tier 2, and
eligible long-term senior unsecured debt held at the holding company level).
This imposes a floor on the long-term unsecured debt stack and directly
interacts with the funding mix optimization: TLAC-eligible issuance is more
expensive than secured funding, but it improves resolution feasibility and
is a hard constraint on the liability optimization problem.
\end{remark}

A bank holding excess capital beyond these buffers faces pressure from
shareholders to deploy it efficiently via buybacks, dividends, or accretive
acquisitions.  In practice, the binding constraint often rotates between
CET1, the leverage ratio (SLR), and TLAC depending on the balance sheet
composition and the interest rate environment.

% ══════════════════════════════════════════════════════════════════════════════
\section{Liquidity Optimization}
% ══════════════════════════════════════════════════════════════════════════════

\subsection{The Liquidity Drag and HQLA Portfolio Management}

LCR and NSFR compliance requires the bank to hold a portfolio of
High-Quality Liquid Assets (HQLA).  HQLA earns below-market yields (e.g.,
Treasuries at 4.5\% vs.\ corporate loans at 6.5\%+), creating a structural
``liquidity drag'' on NIM:

\begin{equation}
  \text{Liquidity\_Drag} = (\bar{r}^A - r^{\text{HQLA}}) \cdot \frac{\text{HQLA}}{A}
\end{equation}

At Apex with HQLA of \$45B and a balance sheet of \$285B (16\% HQLA ratio):

\begin{equation}
  \text{Liquidity\_Drag} \approx (6.2\% - 4.5\%) \times 16\% = 27\, \text{bps NIM}
\end{equation}

HQLA portfolio optimization seeks to maximize yield within the HQLA
composition constraints:

\begin{equation}
  \max_{\text{HQLA mix}} \sum_{i \in \text{HQLA}} w_i \cdot r_i
  \quad \text{s.t.} \quad
  \begin{cases}
    \text{Level 1} \geq 60\% \text{ of HQLA} \\
    \text{Level 2A} \leq 40\% \text{ of HQLA} \\
    \text{Level 2B} \leq 15\% \text{ of HQLA} \\
    \text{Duration} \leq D_{\text{max}} \\
    \text{LCR} \geq 110\% \text{ (internal buffer)}
  \end{cases}
\end{equation}

Within Level 1, the bank can substitute between cash (0\% yield), central bank
reserves (IORB: currently $\approx 4.40\%$ as of April 2026 following the
2024--2025 Fed easing cycle), and Treasury bills/bonds (4.2\%--4.6\% along
the curve).  The optimization shifts toward Treasuries and away from excess cash.
Note that IORB rates and Treasury yields are model parameters that must be
refreshed monthly; the numbers here reflect Q1 2026 conditions.

\subsection{Intraday Liquidity Management}

Intraday liquidity management is distinct from LCR/NSFR.  The bank must fund
settlement obligations throughout the day via Fedwire and CHIPS:

\begin{equation}
  \text{Intraday\_Peak}(t) = \max_{\tau \in [0, T_{\text{day}}]}
  \left[ \text{Outflows}(\tau) - \text{Inflows}(\tau) - \text{Reserve}(\tau) \right]^+
\end{equation}

Minimizing the peak intraday liquidity requirement (and thus the cost of
intraday credit from the Fed) is a treasury optimization problem.  Tools
include:
\begin{itemize}[noitemsep]
  \item Prefunding (posting collateral to Fed discount window)
  \item Payment sequencing (delaying non-urgent payments to afternoon)
  \item Correspondent banking netting
\end{itemize}

\subsection{Funding Mix Optimization}

The liability-side optimization seeks the lowest-cost, most stable funding mix
consistent with NSFR compliance:

\begin{equation}
  \min_{\Delta L_j} \sum_j r_j^L \cdot L_j
  \quad \text{s.t.} \quad
  \text{NSFR} \geq 100\%, \quad \text{LCR} \geq 100\%
\end{equation}

The trade-off: demand deposits are cheap (near-zero rate) but have high NSFR
ASF factors only if stable (50\% core vs.\ 0\% non-stable).  Long-term
unsecured debt has a high ASF (100\%) but is expensive (e.g.,
OIS + 80--120bps for senior unsecured).

\begin{definition}[Optimal Funding Tenor]
Given a flat (or inverted) yield curve, the optimal funding maturity $\tau^*$
for a $T$-year asset satisfies:

\begin{equation}
  \tau^* = \argmin_\tau \left[ r^L(\tau) + \text{Rollover\_Risk}(\tau, T) \right]
\end{equation}

where $\text{Rollover\_Risk}(\tau, T) = \sigma_r^2 \cdot (T - \tau)$ is the
expected cost of refunding if short-term rates rise.
\end{definition}

In a normal upward-sloping curve, short-term funding is cheaper but
rollover risk is higher.  In an inverted curve (2022--2024), long-term
funding may be cheaper even without considering rollover risk.

% ══════════════════════════════════════════════════════════════════════════════
\section{Asset-Liability Management (ALM)}
% ══════════════════════════════════════════════════════════════════════════════

\subsection{NII Sensitivity and the Repricing Gap}

The repricing gap model described in Section~2 provides a
first-order linear approximation of NII sensitivity:

\begin{equation}
  \Delta\text{NII} \approx \text{RSA}^{12m} \cdot \Delta r - \text{RSL}^{12m} \cdot \Delta r
  = \text{Gap}^{12m} \cdot \Delta r
  \label{eq:gap}
\end{equation}

where $\text{RSA}^{12m}$ is the rate-sensitive asset balance repricing within
12 months, and $\text{RSL}^{12m}$ is the corresponding liability balance.

The repricing gap captures only the \emph{volume effect} of rate changes.
Higher-order effects (yield curve non-parallelism, option costs, basis) require
more sophisticated modelling.

\subsection{EVE and Duration Gap}

The Economic Value of Equity sensitivity:

\begin{equation}
  \frac{\Delta\text{EVE}}{\text{EVE}} \approx -D_{\text{gap}} \cdot \frac{\Delta r}{1 + r}
\end{equation}

where the duration gap is:

\begin{equation}
  D_{\text{gap}} = D_A - \frac{L}{A} \cdot D_L
\end{equation}

Positive $D_{\text{gap}}$ means the bank is asset-sensitive in EVE terms
(rising rates reduce EVE).  Negative $D_{\text{gap}}$ means liability-sensitive.

\begin{proposition}[EVE-NII Trade-off]
A bank cannot simultaneously minimize both NII sensitivity and EVE sensitivity
to interest rate risk.  Shortening asset duration increases NII in a rising
rate environment but reduces EVE (less benefit from discounting liabilities
at higher rates).  This is the fundamental ALM trade-off.
\end{proposition}

\begin{remark}[IRRBB Outlier Test and EVE Risk Appetite]
Under BCBS 368 (\textit{Interest Rate Risk in the Banking Book}, 2016) and
its domestic implementations (e.g., the Fed's 2023 IRRBB guidance), supervisors
apply a standardized \emph{outlier test}: a bank whose EVE declines by more
than $15\%$ of Tier 1 Capital under a $\pm 200$\,bp parallel shock is flagged
as an outlier and subject to enhanced supervisory scrutiny or mandatory capital
add-ons.  Formally:

\begin{equation}
  |\Delta\text{EVE}|_{\pm 200\text{bp}} \leq 0.15 \times \text{Tier1}
\end{equation}

At Apex (Tier 1 $\approx \$27B$), this implies an EVE sensitivity limit of
$\approx \$4.0B$ per $\pm 200$\,bp shock.  The risk appetite band $[D^-, D^+]$
on the duration gap should be calibrated to keep the bank comfortably inside
this threshold.  The SVB collapse in March 2023 was in part attributable to
an EVE sensitivity \emph{exceeding 100\% of equity} under a $+300$\,bp scenario ---
a failure of this exact governance guardrail.
\end{remark}

The optimal duration gap maximizes the utility function:

\begin{equation}
  \max_{D_{\text{gap}}} \quad U(\Delta\text{NII}, \Delta\text{EVE})
  = w_1 \cdot \E[\Delta\text{NII}] - w_2 \cdot \text{Var}(\Delta\text{NII})
    - w_3 \cdot |\Delta\text{EVE}|^2
\end{equation}

subject to duration gap within the risk appetite band $[D^-, D^+]$.

\subsection{Non-Maturity Deposit Behavioral Modeling}

Non-maturity deposits (NMDs) — demand deposits, savings accounts, money market
accounts — are contractually repayable on demand but in practice exhibit far
longer behavioural maturities.  Modelling these correctly is the single most
important and most model-risk-intensive assumption in bank ALM.

\paragraph{Deposit Beta.}
The \emph{deposit beta} $\beta_d$ measures how much a deposit rate moves in
response to a 100\,bp change in the policy rate:

\begin{equation}
  \Delta r_{\text{deposit}} = \beta_d \cdot \Delta r_{\text{FF}}
\end{equation}

Empirically, $\beta_d$ varies by product: checking accounts $\approx 0.10$--$0.25$
(highly sticky, slow pass-through), savings accounts $\approx 0.30$--$0.55$,
and retail money market accounts $\approx 0.60$--$0.85$.  A bank with
predominantly checking-account deposits has a substantial NIM advantage in a
rising rate environment: liability costs rise only 10--25\% as fast as
asset yields.

\paragraph{Effective Maturity and Core / Non-Core Classification.}
For ALM and IRRBB purposes, NMDs are modeled with a \emph{behavioural maturity}
that reflects historical runoff patterns:

\begin{equation}
  D_{\text{deposit}}^{\text{effective}} = \sum_{\tau} s_\tau \cdot \tau
\end{equation}

where $s_\tau$ is the fraction of the deposit balance with modelled maturity
$\tau$.  The Basel IRRBB guidelines cap the effective maturity at 5 years for
retail NMDs and 4 years for wholesale NMDs.  Balances are split into:
\begin{itemize}[noitemsep]
  \item \textbf{Core (stable) deposits}: the fraction expected to remain on
        balance sheet even in a stressed environment; ASF factor = 90\%.
        Modelled with longer effective maturity (3--5 years).
  \item \textbf{Non-core (volatile) deposits}: rate-sensitive or
        institutionally-driven balances; ASF factor = 50\%; modelled as
        effectively short-dated (3--12 months).
\end{itemize}

\paragraph{Model Risk Warning.}
NMD behavioral models are among the highest-risk ALM models under SR 11-7.
The 2022--2023 rate cycle exposed significant over-estimation of deposit
stickiness at several mid-sized US banks: depositors migrated to higher-yielding
alternatives (money market funds, T-bills) at a pace not captured in pre-2022
behavioral models.  Any NMD model must include:
\begin{enumerate}[noitemsep]
  \item A stress scenario where deposit beta = 1.0 (full pass-through) to
        test NII sensitivity in a worst-case repricing scenario.
  \item A run-off scenario calibrated to 2023 stress events (not just
        pre-GFC history).
  \item Annual validation comparing model-predicted deposit rates to realized
        rates, with explicit back-testing of the beta assumption.
\end{enumerate}

If the current duration gap exceeds the target, the ALM desk uses interest rate
swaps to close the gap:

\begin{itemize}[noitemsep]
  \item \textbf{Receive-fixed IRS (``receiver swap''):} The bank receives the
        fixed leg and pays floating.  Overlaid on a floating-rate liability,
        this converts the net liability cost to fixed, \emph{increasing}
        effective liability duration and thereby \emph{reducing} the duration
        gap.  Used when the bank is over-positioned asset-sensitive and wishes
        to extend liability duration.
  \item \textbf{Pay-fixed IRS (``payer swap''):} The bank pays the fixed leg
        and receives floating.  Overlaid on a fixed-rate asset (e.g., a
        fixed-rate loan or AFS bond), this converts the net asset return to
        floating, \emph{reducing} effective asset duration.  Used when the
        bank is too asset-sensitive and wants to shorten the asset side of
        the gap.
\end{itemize}

\begin{remark}[Swap Terminology Clarity]
``Payer'' and ``receiver'' always refer to the \emph{fixed leg}: a payer IRS
pays fixed/receives floating; a receiver IRS receives fixed/pays floating.
This convention is universal in inter-dealer markets and must not be confused
with the underlying asset or liability being hedged.
\end{remark}

The notional of swap required to close a duration gap of $\Delta D$ over a
target balance sheet of $A$:

\begin{equation}
  N_{\text{swap}} = \frac{\Delta D \cdot A}{D_{\text{swap}}}
\end{equation}

where $D_{\text{swap}} = \text{DV01}_{\text{swap}} / \text{Price}_{\text{swap}}$.

\subsection{Convexity and Prepayment Risk}

Fixed-rate mortgages and callable bonds exhibit \emph{negative convexity}:
when rates fall, borrowers prepay (shortening asset duration when the bank
least wants it) and when rates rise, borrowers extend (lengthening duration
when the bank least wants that).

The effective duration and convexity:

\begin{equation}
  P(r) \approx P(r_0) \left[ 1 - D_{\text{eff}} \cdot \Delta r
    + \frac{1}{2} C_{\text{eff}} \cdot (\Delta r)^2 \right]
\end{equation}

For a pass-through MBS, $C_{\text{eff}} < 0$, creating a P\&L concavity
that increases with the size of the rate move.

A bank with a large MBS portfolio must hedge convexity separately using:
\begin{itemize}[noitemsep]
  \item Receiver swaptions (provide positive convexity when rates fall)
  \item CMS caps/floors
  \item Callable bond portfolios to offset MBS negative convexity
\end{itemize}

\subsection{AFS vs.\ HTM Classification and AOCI Risk}

The accounting classification of the securities portfolio has direct
consequences for capital management:

\begin{itemize}[noitemsep]
  \item \textbf{Available-for-Sale (AFS):} Marked to market through
        \emph{Other Comprehensive Income} (OCI).  Unrealized losses
        \emph{reduce} AOCI and therefore reduce CET1 capital (for banks
        subject to AOCI opt-out removal or those that have not opted out
        under 12 CFR 3.22(b)).  AFS assets have high HQLA eligibility.
  \item \textbf{Held-to-Maturity (HTM):} Carried at amortized cost; unrealized
        losses do not flow through AOCI or reduce CET1.  However, the intent
        to hold to maturity must be demonstrable --- reclassification ``taints''
        the entire HTM portfolio.  HTM assets are \emph{not} HQLA-eligible.
\end{itemize}

\paragraph{AOCI Capital Risk.}
The 2022--2023 rate rise created a systemic AOCI problem: banks that had
accumulated large AFS securities portfolios during the ZIRP era suffered
mark-to-market losses that eroded CET1 ratios.  The canonical failure:
Silicon Valley Bank held \$91B of AFS/HTM securities against a
\$173B balance sheet; a +200\,bp rate move generated unrealized losses
exceeding its equity.  The integrated CET1 sensitivity to a rate shock:

\begin{equation}
  \Delta\text{CET1} \approx -D_{\text{securities}} \cdot A_{\text{securities}} \cdot \Delta r
                           + \text{NII benefit} \cdot \text{PV multiplier}
\end{equation}

The AOCI effect dominates in the short run; the NII benefit accrues over
multiple years.  This is the \emph{duration trap} reverse stress scenario
identified in Section~8.

\paragraph{Practical Guidance.}
An experienced balance sheet manager will:
\begin{enumerate}[noitemsep]
  \item Monitor the aggregate unrealized gain/loss in the AFS portfolio and
        the implied CET1 impact daily (not just at quarter-end).
  \item Ladder the securities portfolio to limit duration concentration ---
        avoid excessive accumulation of 10-year Treasuries or Agency MBS
        funded by short-dated liabilities.
  \item Consider HTM classification only for assets genuinely intended to be
        held to maturity, with documented liquidity-independent funding to
        back those positions.
  \item Ensure the ALM committee receives a regular report on the ``AOCI
        CET1 bridge'' (unrealized P\&L versus capital buffers) as a standing
        agenda item.
\end{enumerate}

% ══════════════════════════════════════════════════════════════════════════════
\section{Fund Transfer Pricing (FTP)}
% ══════════════════════════════════════════════════════════════════════════════

\subsection{Matched-Maturity FTP Architecture}

The FTP framework transfers the cost of funding and the benefit of funding
between the business lines and a central Treasury pool.  Every asset originates
a charge; every liability originates a credit:

\begin{equation}
  \text{FTP\_Charge}(k) = \underbrace{r^{\text{OIS}}(\tau_k)}_{\text{Funding cost}}
                        + \underbrace{\lambda_k}_{\text{Liquidity premium}}
                        + \underbrace{\psi_k}_{\text{Optionality premium}}
\end{equation}

where:
\begin{itemize}[noitemsep]
  \item $r^{\text{OIS}}(\tau_k)$: OIS swap rate at the contractual maturity $\tau_k$
  \item $\lambda_k$: liquidity premium representing the cost of term funding above OIS
  \item $\psi_k$: optionality premium for embedded options (prepayment, early
        withdrawal, commitment fee)
\end{itemize}

\subsection{FTP and the Elimination of Cross-Subsidy}

Without FTP, a bank using short-term cheap deposits to fund long-term loans
appears to generate a large NII spread.  In reality, it is accumulating
unhedged rollover risk.  FTP corrects this:

\begin{equation}
  \text{NIM}_{\text{adjusted}}(k) = r_k^A - \text{FTP}(k) - \text{EL}_k
\end{equation}

A loan desk with high gross NIM but high FTP charge (long-dated fixed-rate)
may be less attractive than a shorter-duration product with lower gross NIM
but a lower FTP charge.

\subsection{Liquidity Premium Components}

The liquidity premium $\lambda$ contains:

\begin{equation}
  \lambda = \lambda_{\text{term}} + \lambda_{\text{LCR}} + \lambda_{\text{NSFR}} + \lambda_{\text{CFP}}
\end{equation}

\begin{itemize}[noitemsep]
  \item $\lambda_{\text{term}}$: term funding spread above OIS (market-observable from
        senior unsecured issuance)
  \item $\lambda_{\text{LCR}}$: cost of holding HQLA attributable to this liability class
        (high-runoff liabilities incur higher LCR charges)
  \item $\lambda_{\text{NSFR}}$: cost of providing stable funding to match RSF on the
        asset side
  \item $\lambda_{\text{CFP}}$: cost of maintaining Contingency Funding Plan reserves
\end{itemize}

\subsection{Governance and Incentive Design}

FTP design creates strong incentives.  If FTP is set too low:
\begin{itemize}[noitemsep]
  \item Loan desks over-originate because the cost of funding is underpriced
  \item Balance sheet grows unsustainably
\end{itemize}

If FTP is set too high:
\begin{itemize}[noitemsep]
  \item Desks under-price deposits relative to true value
  \item Treasury retains excessive profit; business lines cross-subsidize Treasury
\end{itemize}

FTP rates must be reviewed at minimum quarterly by the ALM Committee,
with Treasury as the internal market-maker.

% ══════════════════════════════════════════════════════════════════════════════
\section{RAROC and Risk-Adjusted Pricing}
% ══════════════════════════════════════════════════════════════════════════════

\subsection{RAROC Framework}

Risk-Adjusted Return on Capital (RAROC) is the primary metric for comparing
across business lines with different risk profiles:

\begin{equation}
  \text{RAROC}(k) =
  \frac{(r_k^A - r_k^{\text{FTP}}) \cdot A_k - \text{EL}_k \cdot A_k - \text{OpCost}_k}
       {\text{EC}_k}
\end{equation}

\begin{definition}[RAROC Hurdle Rate]
The hurdle rate $h$ is typically the bank's cost of equity, estimated via CAPM:
\begin{equation}
  h = r_f + \beta \cdot (r_m - r_f) = 4.5\% + 1.2 \times 5.5\% = 11.1\%
\end{equation}
for a money-center bank with $\beta \approx 1.2$.  Any business line with
$\text{RAROC} < h$ destroys shareholder value and should be restructured or exited.
\end{definition}

\subsection{Loan Pricing Using RAROC}

The minimum loan rate $r^*$ for a new loan that achieves the hurdle RAROC:

\begin{equation}
  r^* = r^{\text{FTP}}(\tau) + \lambda^{\text{FTP}} + \text{EL} + \frac{h \cdot \text{EC}}{A}
      + \text{OpCost}
\end{equation}

Example: 5-year commercial loan

\begin{align}
  r^{\text{OIS}}(5Y) &= 4.40\% \\
  \lambda^{\text{FTP}}(5Y) &= 0.45\% \\
  \text{EL} = \text{PD} \cdot \text{LGD} &= 1.5\% \times 45\% = 0.675\% \\
  \frac{h \cdot \text{EC}}{A} &= 11.1\% \times 8\% = 0.888\% \\
  \text{OpCost} &= 0.20\% \\
  \hline
  r^* &\approx 6.61\%
\end{align}

Any loan priced above 6.61\% is value-creating.  Below this rate, the bank is
lending at a negative RAROC.

\subsection{Capital Allocation and Performance Measurement}

Economic capital is allocated to business lines proportional to their risk
contribution.  The allocation must satisfy:

\begin{equation}
  \sum_k \text{EC}_k = \text{EC}_{\text{portfolio}}
\end{equation}

Using Euler's theorem on homogeneous risk measures (VaR is 1-homogeneous):

\begin{equation}
  \text{EC}_k = w_k \cdot \frac{\partial \text{VaR}_\alpha}{\partial w_k}
\end{equation}

where $w_k$ is the exposure weight of business line $k$.  This ensures the
diversification benefit is properly distributed.

% ══════════════════════════════════════════════════════════════════════════════
\section{Stress Testing Integration}
% ══════════════════════════════════════════════════════════════════════════════

\subsection{From DFAST to Balance Sheet Strategy}

DFAST stress testing is not merely a regulatory exercise — it is a strategic
planning tool.  The severely adverse CET1 trough determines:

\begin{itemize}[noitemsep]
  \item \textbf{Maximum loan growth capacity}: if the bank originates $\Delta A$
        in new loans before DFAST, the stressed loss on those loans reduces the
        CET1 trough by approximately $\text{EL}_{\text{stress}} \cdot \Delta A / \text{RWA}$.
  \item \textbf{Maximum dividend capacity}: any dividend paid reduces starting
        CET1 and directly lowers the trough.
  \item \textbf{Buyback capacity}: capital returned via buybacks above the buffer
        could cause the trough to breach the internal floor.
\end{itemize}

The balance sheet optimization explicitly incorporates the DFAST constraint
(equation~\eqref{eq:dfast}) as a planning constraint.

\subsection{Reverse Stress Testing}

Reverse stress testing identifies the balance sheet composition that would
cause the bank to breach regulatory minima.  The reverse stress problem:

\begin{equation}
  \min_{\text{scenario}} \text{CET1\_trough}(\text{scenario})
  \quad \text{s.t.} \quad \text{Scenario plausible under risk model}
\end{equation}

Key reverse stress scenarios:

\begin{enumerate}[noitemsep]
  \item \textbf{Credit concentration}: a sector-concentrated credit shock
        (e.g., CRE collapse of -55\%) combined with deposit run causes simultaneous
        asset loss and liability instability.
  \item \textbf{Duration trap / AOCI spiral}: a sharp rate rise (+300bps
        rapidly) reduces both EVE and AFS unrealized values; the AFS mark
        flows through AOCI, directly eroding CET1.  If CET1 approaches the
        SCB-adjusted floor, management credibility is damaged and deposit
        outflows accelerate (the exact SVB 2023 failure mechanism).  The
        key insight: the duration trap is a \emph{solvency-then-liquidity}
        event, not purely a liquidity event.  Reverse stress must therefore
        model the second-round deposit outflow triggered by a visible capital
        deterioration.
  \item \textbf{Liquidity-solvency spiral}: an LCR breach forces fire-sale
        of HQLA, which depresses prices of the HQLA assets themselves (pro-cyclical).
\end{enumerate}

Each reverse stress scenario generates a specific ``distance to failure''
metric:

\begin{equation}
  D_{\text{failure}} = \frac{\text{CET1\_current} - \text{CET1\_at\_failure}}
                            {\text{CET1\_1-year-vol}}
\end{equation}

% ══════════════════════════════════════════════════════════════════════════════
\section{Technology, Calculators, and Infrastructure}
% ══════════════════════════════════════════════════════════════════════════════

\subsection{Required Calculator Stack}

Balance sheet optimization requires an integrated stack of quantitative
engines.  The minimum production set:

\begin{table}[H]
  \centering
  \caption{Core Calculator Stack for Balance Sheet Optimization}
  \begin{tabular}{lll}
    \toprule
    Calculator & Purpose & Key Output \\
    \midrule
    RWA / Capital Engine & SA and IMA RWA; CET1/Tier1/Total ratios & Capital ratios \\
    VaR / SVaR Engine & IMA market risk capital & 10-day 99\% VaR, SVaR \\
    FRTB-SA Engine & SA sensitivities; curvature; DRC & FRTB-SA capital \\
    IFRS 9 ECL Engine & Loan loss provision; staging & ECL coverage ratio \\
    DFAST Engine & 9-quarter CET1 projection & CET1 trough \\
    ALM Engine & NII/EVE sensitivity; repricing gap & Duration gap; NIM \\
    FTP Engine & Matched-maturity funding charges & Desk P\&L (FTP-adjusted) \\
    Collateral / SIMM & VM/IM; liquidity stress & IM requirement; stress call \\
    XVA Suite & CVA/DVA/FVA/MVA/ColVA/KVA & XVA charge per trade \\
    RAROC Calculator & Risk-adjusted pricing; hurdle test & Minimum loan rate \\
    LCR / NSFR Calculator & Liquidity ratio compliance & LCR, NSFR ratios \\
    Concentration Monitor & Single-name/sector/geo limits & Breach flags, HHI \\
    \bottomrule
  \end{tabular}
\end{table}

\subsection{Data Architecture}

The calculators require a consistent, real-time data layer:

\begin{itemize}[noitemsep]
  \item \textbf{Positions layer}: real-time feed from OMS and loan systems;
        all instruments mapped via instrument master (ISIN/CUSIP/product type)
  \item \textbf{Market data layer}: live prices, rates, vol surfaces, credit
        spreads; GBM simulation for real-time VaR (Apex uses 500ms tick interval)
  \item \textbf{Reference data}: CSA terms, counterparty credit ratings, FATF
        grey-list, PEP database, ISDA SIMM risk weights
  \item \textbf{Event log}: append-only audit trail of all balance sheet changes,
        capital actions, and model outputs (SQLite in Apex implementation)
  \item \textbf{Regulatory parameters}: Fed-published DFAST scenarios, Basel
        risk weights, LCR runoff rates (updated annually or upon regulatory change)
\end{itemize}

\subsection{Computational Requirements}

\begin{table}[H]
  \centering
  \caption{Computational Requirements by Calculator}
  \begin{tabular}{lllll}
    \toprule
    Calculator & Frequency & Latency & Compute & Memory \\
    \midrule
    VaR MC (2,000 paths, 11 instruments) & Daily post-close & 50ms & 1 CPU & 50MB \\
    FRTB-SA (delta/vega/curvature) & Daily & 200ms & 4 CPU & 200MB \\
    DFAST (9Q, 3 scenarios) & Quarterly & 5--30s & 8 CPU & 1GB \\
    XVA (1,000 paths, 50 trades) & On trade & 2--10s & 16 CPU & 4GB \\
    SIMM (CRIF-based, all CSAs) & Daily & 100ms & 2 CPU & 100MB \\
    LCR / NSFR & Daily & 500ms & 2 CPU & 200MB \\
    ALM / NII sensitivity & Monthly & 1s & 2 CPU & 100MB \\
    IFRS 9 ECL (50 obligors) & Quarterly & 5ms & 1 CPU & 20MB \\
    \bottomrule
  \end{tabular}
\end{table}

For production-scale balance sheet optimization (850+ obligors, 10,000+
derivative trades, full FRTB-IMA), the compute requirements scale to:
\begin{itemize}[noitemsep]
  \item XVA: 64--128 GPU cores for real-time Monte Carlo (AAD sensitivities)
  \item FRTB IMA: historical simulation on 10-year daily data for 500+ risk factors
  \item DFAST: cluster parallelization across 3 scenarios × 9 quarters
\end{itemize}

\subsection{Dashboards and Reporting}

A best-practice balance sheet optimization dashboard suite consists of:

\paragraph{1.\ Capital Dashboard} CET1/Tier1/Total/Leverage ratios vs.\ regulatory
minimums and internal targets.  RWA decomposition by risk type (credit,
market, operational).  Capital headroom to constraint breaches.  Trajectory
over prior 8 quarters.

\paragraph{2.\ Liquidity Dashboard} LCR and NSFR ratios real-time (intraday
updates for LCR).  HQLA composition by Level (1/2A/2B).  30-day net cash
outflow waterfall.  Intraday liquidity utilization vs.\ credit line.

\paragraph{3.\ ALM / Rate Risk Dashboard} NII sensitivity bar chart
($\pm$100bps, $\pm$200bps, instantaneous).  EVE sensitivity.  Repricing gap
schedule (7 buckets, asset vs.\ liability bars).  Duration gap trend vs.\
risk appetite band.  Warning indicators (SVB-style flag if EVE sensitivity
exceeds --15\%).

\paragraph{4.\ DFAST / Stress Dashboard} 9-quarter CET1 trajectory under all
3 scenarios (Plotly chart with min CET1 floor and internal risk appetite band).
Loss attribution waterfall (credit, market, operational, PPNR).  Distance to
regulatory minimum.

\paragraph{5.\ FTP / RAROC Dashboard} Desk-level FTP-adjusted P\&L ranking.
RAROC vs.\ hurdle rate by business line.  New origination pricing heat-map.
FTP curve visualization (OIS + liquidity premium by tenor).

\paragraph{6.\ XVA Dashboard} CVA, DVA, FVA, MVA, ColVA, KVA by counterparty
and trade.  Live WebSocket updates on new trade submissions.  XVA P\&L explain
(delta, vega, credit) vs.\ prior day.

\paragraph{7.\ Boardroom / Scenario Dashboard} Real-time agent-driven scenario
narrative.  Scenario engine with live risk factor injection.  Model registry
and MDD governance status.

% ══════════════════════════════════════════════════════════════════════════════
\section{How Balance Sheet Optimization Generates Revenue}
% ══════════════════════════════════════════════════════════════════════════════

\subsection{Revenue Sources by Lever}

Balance sheet optimization generates value through several distinct mechanisms:

\subsubsection{Capital Efficiency: Releasing Trapped Capital}

When RWA can be reduced without sacrificing income, the freed capital can be
redeployed into higher-RAROC activities or returned to shareholders.

\textbf{Example -- Loan portfolio restructuring:}
A bank holds \$10B of low-grade CRE loans (RW = 100\%, generating $6\%$ yield).
These can be replaced with investment-grade corporate bonds (RW = 20\%,
yielding $5.5\%$):

\begin{align}
  \text{RWA reduction} &= (100\% - 20\%) \times \$10B = \$8B \\
  \text{CET1 freed} &= 7\% \times \$8B = \$560M \\
  \text{NII sacrifice} &= (6\% - 5.5\%) \times \$10B = \$50M/\text{yr} \\
  \text{Return on freed capital} &= \$50M / \$560M = 8.9\% \text{ cost}
\end{align}

If the freed CET1 can be deployed into a business generating RAROC $> 8.9\%$,
the swap is value-creating.

\subsubsection{NII Optimization: Positioning the Duration Gap}

A bank that correctly positions its duration gap before a rate cycle captures
NII uplift:

\textbf{Example -- Rate rise scenario:}
Starting NII \$40B/year, duration gap +0.41yr.
Rate rises 200bps: $\Delta\text{NII} = +\$6.8B/\text{yr}$ (Apex ALM model).
This is a structural P\&L benefit worth capturing by avoiding over-hedging.
A bank with zero duration gap would earn \$0 uplift; correctly calibrating
the gap to +0.41yr (consistent with risk appetite) generates \$6.8B incremental
NII under a +200bps scenario.

Conversely, if rates had fallen 200bps, the bank would have lost \$5.9B NII.
The asymmetric positioning (positive gap in an asset-sensitive bank) reflects
the rate environment view embedded in the ALM strategy.

\subsubsection{Liquidity Premium Capture}

Treasury can earn the term liquidity premium by issuing short-term paper and
lending to the asset book at FTP rates that include the term premium.  This
``carry'' between short funding and long assets is a persistent revenue source
in a normally-sloped yield curve:

\begin{equation}
  \text{Carry} = r^{\text{FTP}}(10Y) - r^{\text{FTP}}(3M)
              = 4.55\% - 5.15\% = -60\text{bps}
\end{equation}

In the current inverted curve, the carry is negative, making it advantageous
to issue long-term and hold short-dated assets — the opposite of the normal regime.

\subsubsection{XVA Optimization: KVA and FVA Pricing}

The bank can generate revenue by correctly pricing XVA charges into derivatives
trades at inception.  If the XVA charge is not priced, the bank delivers P\&L
upfront but accumulates a long-dated liability as the XVA is realized.

\textbf{KVA (Capital Valuation Adjustment):} The cost of holding regulatory
capital against the trade over its lifetime:

\begin{equation}
  \text{KVA} = h \cdot \int_0^T \E[\text{EC}(t)] \cdot \text{df}(t)\, dt
\end{equation}

where $h$ is the hurdle rate and $\text{EC}(t)$ is expected economic capital
at time $t$.  KVA is typically 5--15bps on notional for a 5-year IRS.

Pricing KVA into derivatives trades means the desk charges the client the full
cost of the trade, including the capital cost — generating an immediate
revenue recognition that compensates for the long-dated capital consumption.

\subsubsection{Regulatory Arbitrage: Cleared vs.\ Bilateral}

Clearing eligible derivatives via a CCP (LCH, CME) attracts lower SA-CCR and
IM requirements:

\begin{align}
  \text{IM (bilateral SIMM)} &\approx \$70M \\
  \text{IM (LCH cleared)} &\approx \$25M \\
  \text{IM savings} &= \$45M
\end{align}

At a funding cost of 5.25\%, the IM saving generates an annual FVA/MVA saving
of $\$45M \times 5.25\% = \$2.4M/\text{yr}$ per portfolio of this size.
For a large derivatives book (\$500B notional), clearing-driven IM reduction
can generate \$50--150M/yr in FVA/MVA savings.

\subsubsection{Securitization and Synthetic Risk Transfer}

Loan securitization (ABS, CLO, RMBS, CMBS) and synthetic risk transfer (SRT)
transactions are powerful balance sheet optimization tools that reduce both
capital consumption and balance sheet size:

\begin{itemize}[noitemsep]
  \item \textbf{Traditional securitization (true-sale):} Loans are transferred
        to a special purpose vehicle (SPV); if derecognition criteria are met
        (SFAS 166/167, ASC 860), assets come off balance sheet, reducing both
        RWA and total leverage exposure.  The bank retains a first-loss piece
        (``retained interest'') which absorbs credit risk, but the senior
        tranches are placed with third-party investors.
  \item \textbf{Synthetic SRT:} The bank retains the loans on balance sheet but
        purchases credit protection (CDS or funded guarantee) on the mezzanine
        tranche.  Under Basel III SEC-SA or the SEC-IRBA framework, this can
        reduce RWA by 40--70\% on the reference portfolio.
\end{itemize}

\textbf{Example --- CLO securitization of \$5B leveraged loan portfolio:}
\begin{align}
  \text{Pre-securitization RWA} &= 100\% \times \$5B = \$5B \\
  \text{Retained first-loss (5\%)} &= \$250M \text{ at 1250\% RWA} = \$3.1B \text{ RWA} \\
  \text{RWA reduction} &= \$5B - \$3.1B = \$1.9B \\
  \text{CET1 freed} &= 7\% \times \$1.9B = \$133M
\end{align}

The freed CET1, deployed at RAROC 15\%, generates \$20M/yr of incremental
returns.  SRT transactions typically cost 3--5\% of the protected tranche
notional in guarantee fees, so the economics must be evaluated on a net basis.

\subsubsection{HQLA Yield Optimization}

Within the LCR/NSFR constraints, optimizing the HQLA portfolio composition
generates incremental yield:

\begin{equation}
  \text{HQLA yield gain} = (r^{\text{optimal}} - r^{\text{cash}}) \times \text{HQLA balance}
                         = (4.65\% - 0\%) \times \$45B = \$2.1B/\text{yr}
\end{equation}

Replacing non-earning excess reserves with a duration-matched portfolio of
Treasuries and Agency bonds earns 4.5--4.7\% on \$45B = $\sim$\$2.1B/yr with
no reduction in LCR compliance.

\subsection{Summary: Revenue Impact of Balance Sheet Optimization}

\begin{table}[H]
  \centering
  \caption{Estimated Annual Revenue Impact (Apex Scale, \$285B Balance Sheet)}
  \begin{tabular}{lrr}
    \toprule
    Optimization Lever & Annual P\&L Impact & Key Model / Engine \\
    \midrule
    HQLA yield optimization & +\$2.1B & ALM Engine \\
    NII duration positioning (+200bps scenario) & +\$6.8B (scenario-dep.) & ALM Engine \\
    Clearing-driven IM reduction (MVA/FVA saving) & +\$50--150M & SIMM / XVA \\
    KVA pricing into new derivatives & +\$100--250M & XVA Suite \\
    CET1 release via lower-RWA assets & +\$200--500M (redeployed) & Capital Engine \\
    Securitization / SRT (CET1 redeployment) & +\$100--300M & Capital Engine \\
    AOCI management (AFS/HTM laddering) & Avoid \$1--3B AOCI capital erosion & ALM Engine \\
    Loan pricing discipline (RAROC floor) & Prevents erosion of 50--100bps NIM & RAROC \\
    FTP governance (correct term premium) & Eliminates \$0.5--1.5B cross-subsidy & FTP Engine \\
    \bottomrule
  \end{tabular}
\end{table}

The cumulative revenue uplift from a mature balance sheet optimization program
relative to an unoptimized baseline is typically 15--40bps on total assets, or
\$0.4--1.1B/yr at Apex scale.  At JPMorgan scale (\$3.9T balance sheet),
the program generates \$3--8B/yr in incremental returns.

% ══════════════════════════════════════════════════════════════════════════════
\section{Governance and Model Risk}
% ══════════════════════════════════════════════════════════════════════════════

\subsection{SR 11-7 Governance for Optimization Models}

Every quantitative model used in balance sheet optimization is subject to
SR 11-7 model risk governance:

\begin{itemize}[noitemsep]
  \item \textbf{Model Development Document (MDD):} Each calculator (VaR, DFAST,
        ALM, FTP, SIMM, XVA) requires a complete MDD including theoretical
        basis, implementation, validation, limitations, use authorization, and
        open findings.
  \item \textbf{Independent validation:} All Tier 1 models must be validated
        by a Model Validation Officer independent of the model developer.
  \item \textbf{Ongoing monitoring:} Model KRIs (key risk indicators) tracked
        monthly: VaR backtesting exception rate, SIMM vs.\ actual IM
        reconciliation, DFAST vs.\ actual CET1 variance, NII model vs.\ actual
        NII error.
  \item \textbf{Use authorization:} Each model's Section 7 specifies authorized
        uses, prohibited uses, and the approval chain.
\end{itemize}

\subsection{ALM Committee Governance}

The ALM Committee is the governing body for balance sheet optimization
decisions.  Membership:

\begin{itemize}[noitemsep]
  \item CRO (Chair)
  \item CFO / Treasurer
  \item Head of Trading (business line representative)
  \item Head of Commercial Banking
  \item Head of Model Validation (advisory)
  \item Chief Risk Economist (macro forecasting)
\end{itemize}

The Committee meets monthly and reviews: NII sensitivity, EVE sensitivity,
duration gap vs.\ target, LCR/NSFR ratios, FTP rates, and any proposed
balance sheet transactions above materiality thresholds.

\subsection{The Three Lines of Defence}

\begin{enumerate}[noitemsep]
  \item \textbf{First line (Treasury / business lines):} Runs the ALM engine,
        FTP, and RAROC calculations.  Originates hedging recommendations.
  \item \textbf{Second line (Risk / Model Validation):} Independent review of
        ALM models and FTP rates.  Validates DFAST projections.  Monitors
        model KRIs.
  \item \textbf{Third line (Internal Audit):} Annual audit of the balance sheet
        optimization program; testing of governance controls, data integrity,
        and regulatory compliance.
\end{enumerate}

% ══════════════════════════════════════════════════════════════════════════════
\section{Conclusion}
% ══════════════════════════════════════════════════════════════════════════════

Balance sheet optimization is the highest-leverage quantitative discipline
in banking.  Its scope encompasses capital structure, liquidity management,
interest rate risk, pricing discipline, and regulatory strategy.  It is not
a single model or a single optimization problem --- it is an integrated system
of models, calculators, governance processes, and strategic decisions that
operates continuously over the life of the bank.

The key insight is that the regulatory constraints (capital ratios, LCR, NSFR,
SCB, TLAC) are not merely compliance requirements --- they are \emph{priced inputs}
to the optimization.  A bank that treats them as binding constraints to be
minimally satisfied leaves significant value on the table.  A bank that treats
them as the efficient frontier of the feasible set, and optimizes within it,
can generate several hundred basis points of incremental RORWA relative to
peers.

\paragraph{Lessons from the 2022--2023 Cycle.}
The rate rise of 2022--2023 and the subsequent regional bank failures of 2023
provided a stark reminder of the interconnection between the balance sheet
optimization dimensions:
\begin{itemize}[noitemsep]
  \item \textbf{AOCI is capital}: unrealized losses in the securities portfolio
        are real solvency risk, not just accounting noise.
  \item \textbf{Deposit beta assumptions are fragile}: behavioral models built
        on pre-2015 data underestimated the speed of deposit migration in a
        digital banking environment.
  \item \textbf{Duration mismatch is fatal at the extremes}: the NIM benefit
        of running a positive duration gap does not compensate for the
        existential risk of an AOCI-driven CET1 breach.
  \item \textbf{Concentration in the liability mix amplifies failure}: banks
        with concentrated, rate-sensitive, non-insured depositor bases are
        inherently more fragile than those with a diversified, sticky retail
        deposit franchise.
\end{itemize}

\paragraph{Open Priorities.}
At Apex Global Bank, the near-term balance sheet optimization priorities are:
\begin{enumerate}[noitemsep]
  \item Wire \texttt{CollateralEngine} balances into the CVA exposure
        calculation (TODO-025) to ensure XVA reflects post-collateral exposure.
  \item Build a deposit beta dashboard that monitors realized vs.\ model beta
        monthly and triggers an ALM committee review if the spread exceeds 15\,bp.
  \item Implement AOCI sensitivity reporting into the capital dashboard with
        daily granularity, including a ``capital at risk from AOCI'' metric.
  \item Evaluate the TLAC eligible debt stack and model the issuance cost
        trade-off between callable senior preferred and bail-in-able Tier 2.
\end{enumerate}

\paragraph{Dynamic Optimization: The Next Frontier.}
The next frontier is dynamic balance sheet optimization: replacing the current
static quarterly cycle with a continuous, model-driven optimization that
updates recommendations in real-time as market conditions change.  This requires
AAD (Adjoint Algorithmic Differentiation) sensitivities across all models,
GPU-accelerated Monte Carlo, and reinforcement learning for the hedging
decision layer --- capabilities that are commercially available and being
deployed at G-SIBs today.  The Apex infrastructure is positioned to support
this transition, with the real-time market data feed, event-log audit trail,
and WebSocket-enabled dashboards providing the data backbone for a real-time
optimization loop.

% ── Bibliography ──────────────────────────────────────────────────────────────
\newpage
\begin{thebibliography}{99}
\bibitem{basel3}
  Basel Committee on Banking Supervision (2017).
  \textit{Basel III: Finalising post-crisis reforms}.
  Bank for International Settlements.

\bibitem{frtb}
  Basel Committee on Banking Supervision (2016).
  \textit{Minimum capital requirements for market risk (FRTB)}, d352.
  Bank for International Settlements.

\bibitem{sr117}
  Federal Reserve (2011).
  \textit{Supervisory Guidance on Model Risk Management}, SR 11-7.

\bibitem{sr1217}
  Federal Reserve (2012).
  \textit{Supervisory Guidance on Stress Testing for Banking Organizations with More than \$10 Billion in Total Consolidated Assets}, SR 12-7.

\bibitem{dfast}
  Federal Reserve (2022).
  \textit{Dodd-Frank Act Stress Tests: Final Rule}, 12 CFR 252.

\bibitem{lcr}
  Basel Committee on Banking Supervision (2013).
  \textit{Basel III: The Liquidity Coverage Ratio and liquidity risk monitoring tools}.

\bibitem{nsfr}
  Basel Committee on Banking Supervision (2014).
  \textit{Basel III: The Net Stable Funding Ratio}.

\bibitem{simm}
  ISDA (2023).
  \textit{ISDA SIMM Methodology, Version 2.6}.
  International Swaps and Derivatives Association.

\bibitem{hull}
  Hull, J.\ C.\ (2022).
  \textit{Options, Futures, and Other Derivatives}, 11th ed.
  Pearson.

\bibitem{fabozzi}
  Fabozzi, F., Mann, S.\ V.\ (2011).
  \textit{The Handbook of Fixed Income Securities}, 8th ed.
  McGraw-Hill.

\bibitem{saunders}
  Saunders, A., Cornett, M.\ M.\ (2018).
  \textit{Financial Institutions Management: A Risk Management Approach}, 9th ed.
  McGraw-Hill.

\bibitem{choudhry}
  Choudhry, M.\ (2012).
  \textit{The Principles of Banking}.
  Wiley.

\bibitem{iwanowski}
  Iwanowski, R., Harrington, B.\ (2007).
  ``Balance Sheet Optimization for Financial Institutions,''
  \textit{Journal of Applied Finance}, 17(1), 28--45.
\end{thebibliography}

\end{document}
